\subsection{Experimental Methodology}

We evaluated the system using a controlled demo with both benchtop and short-range field testing. The setup
used four LoRa nodes (Node A to Node D), all running the same firmware as equal peers in the mesh (no
coordinator/master node). For repeatability, nodes were given fixed GPS coordinates (either physically measured
or hardware-mocked) so the spatial geometry stayed consistent while testing individual subsystems. Across all
test cases, Node A acted as the message originator, Nodes B and C were intermediate relay candidates, and Node D
was the terminal receiver. We used separate monitoring terminals to collect relay timestamps and packet logs
from all nodes, giving us a common time reference for cross-node latency measurements. All LoRa transceivers
operated at 433~MHz with a spreading factor of SF9, a bandwidth of 125~kHz, a coding rate of 4/5, and a nominal
transmit power of +17~dBm, yielding a theoretical sensitivity floor of approximately $-137$~dBm and an estimated
line-of-sight range of 2~km per hop.

\subsection{Multi-Hop Mesh Routing}

To test the multi-hop routing capability, Nodes A and D were placed beyond the single-hop range possible with
the configured transmit power, so there was no direct radio path between them. Node B was placed in between,
within range of both ends. Node A transmitted a test message with a maximum hop count of three. Node B received
the packet at an RSSI of $-108$~dBm and an SNR of $-3$~dB, delivered it to its local application layer,
decremented the hop count, and relayed it after its computed hold-back delay expired. Node D then received the
relayed packet at an RSSI of $-113$~dBm and an SNR of $-5$~dB, confirmed local delivery, and the relay chain
ended since no other eligible candidates were within range. End-to-end delivery finished in under one second,
and the hop-count field was confirmed to decrement correctly at each hop. This verifies that the mesh routing
protocol can extend communication range beyond a single node's direct transmission range, without using
pre-configured routing tables or infrastructure.

\subsection{Priority-Based Delivery}

To assess QoS enforcement, we injected two messages from Node A in quick succession: a LOW priority message
followed by an EMERGENCY priority message, separated by about 150~ms. Both were routed through Node B toward
Node C. From Node B's relay logs, the hold-back delay was about 320~ms for the LOW priority message and about
75~ms for the EMERGENCY message. Delivery timestamps at Node C confirmed that the EMERGENCY message arrived
before the LOW priority message even though it was transmitted later. The observed ratio of hold-back delays
(approximately $4.3\times$) matches the intended fourfold priority scaling in the delay formula. This confirms
that the QoS classification carries across at least one relay hop and still works when different priority
messages are queued at the same intermediate node.

\subsection{Directional Propagation}

We evaluated geographic propagation control using a four-node setup where Nodes B and C were within a
northward angular arc of approximately $\pm 40^\circ$ relative to Node A, while Node D was at a bearing of
approximately $120^\circ$ east of north, clearly outside a $\pm 45^\circ$ cone pointing north. In OMNI mode, a
broadcast message with a maximum hop count of two was received and relayed by all three surrounding nodes,
resulting in four total relay events including the origination transmission. We then repeated the test in
DIRECTIONAL mode with a target heading of $0^\circ$ (due north), which defines a $\pm 45^\circ$ propagation cone.
Nodes B and C (inside the cone) received and relayed as expected. Node D still received and delivered the
packet locally (so radio reception worked), but its network layer computed the bearing from Node A to Node D as
$120^\circ$ from north, saw that it exceeds the cone boundary, and suppressed re-transmission using the
deliver-only policy. Switching from OMNI to DIRECTIONAL mode reduced total relay events by 33\% for this
topology. The filtering was fully autonomous at each node using independent Haversine-based bearing
computation from the coordinates in the packet header, with no centralized arbitration.

\subsection{Intelligent Relay Timing}

We evaluated the hold-back mechanism by placing three relay candidates (Nodes B, C, and D) at distances of
approximately 200~m, 900~m, and 1.6~km from Node A. The received signal conditions were: Node B at $-85$~dBm
RSSI and $+6$~dB SNR, Node C at $-105$~dBm RSSI and $-1$~dB SNR, and Node D at $-118$~dBm RSSI and $-7$~dB
SNR. Node A broadcast a single MEDIUM priority message and all three nodes received it at the same time. Each
node computed its hold-back delay independently using its locally measured signal quality. Node D (farthest and
weakest signal) relayed first at about 82~ms after origination. Node C relayed at about 194~ms. Node B (closest
and strongest signal) computed a delay of about 340~ms, but before its timer expired it overheard Node C's
re-transmission of the same packet (matching originator identifier and sequence number) and canceled its own
pending relay. So we got two relay transmissions from three candidates, which is a 33\% reduction in redundant
channel occupancy for this configuration. The ordering we saw (farthest node first, closest node suppressed) is
the intended emergent behavior from the signal-quality-weighted delay formula. This confirms distributed
best-relay selection without any explicit inter-node negotiation.

\subsection{Summary}

Across the evaluated subsystems, the field tests confirmed four key behaviors: multi-hop delivery beyond direct
range, priority-consistent forwarding delay, directional relay suppression, and distributed best-relay
selection. We did not observe routing loops, duplicate final deliveries, or spurious relay activations during
the evaluation period. These results validate decentralized operation for the tested topology and radio
settings, but broader generalization would need larger-scale mobility trials.